\documentclass[11pt,a4paper]{article}

% ----------------------------------------------------------------------------
% Packages
% ----------------------------------------------------------------------------
\usepackage[margin=1in]{geometry}
\usepackage{amsmath, amssymb, amsthm}
\usepackage{microtype}
\usepackage{xcolor}
\usepackage[colorlinks=true, linkcolor=blue!50!black, citecolor=blue!50!black,
            urlcolor=blue!50!black]{hyperref}
\usepackage[ruled,vlined,linesnumbered]{algorithm2e}
\usepackage{booktabs}
\usepackage{enumitem}

% ----------------------------------------------------------------------------
% Theorem environments
% ----------------------------------------------------------------------------
\newtheorem{theorem}{Theorem}[section]
\newtheorem{lemma}[theorem]{Lemma}
\newtheorem{proposition}[theorem]{Proposition}
\newtheorem{corollary}[theorem]{Corollary}

\theoremstyle{definition}
\newtheorem{definition}[theorem]{Definition}
\newtheorem{construction}[theorem]{Construction}

\theoremstyle{remark}
\newtheorem{remark}[theorem]{Remark}

% ----------------------------------------------------------------------------
% Custom macros — domain-specific
% ----------------------------------------------------------------------------
% Fields
\newcommand{\binF}{\mathbb{F}_{2^{128}}}        % Binius oracle field (B128/GHash)
\newcommand{\akF}{\mathbb{F}_{p}}                % Akita prime field (fp128)
\newcommand{\F}{\mathbb{F}}                       % generic field
\newcommand{\GF}[1]{\mathbb{F}_{#1}}             % \GF{2}, \GF{2^k}

% Common objects
\newcommand{\BiniusScalar}{\mathsf{BiniusScalar}}
\newcommand{\AkitaFieldScalar}{\mathsf{AkitaFieldScalar}}
\newcommand{\eq}{\mathrm{eq}}                    % equality polynomial
\newcommand{\MLE}[1]{\widetilde{#1}}              % multilinear extension overlay

% Protocol-specific
\newcommand{\BitTable}{B}                         % the bit-table multilinear
\newcommand{\Transparent}{t}                      % transparent oracle vector
\newcommand{\Message}{m}                          % message / committed oracle vector
\newcommand{\selectedPoint}{r_s}                  % output of selected-sum sumcheck
\newcommand{\boolPoint}{r_b}                      % output of booleanity sumcheck
\newcommand{\reducedPoint}{r_*}                   % output of claim-reduction sumcheck
\newcommand{\evalSelected}{e_s}                   % B(\selectedPoint)
\newcommand{\evalBool}{e_b}                       % B(\boolPoint)

% Crypto shorthand
\newcommand{\bits}{\{0,1\}}
\newcommand{\negl}{\mathsf{negl}}
\DeclareMathOperator{\Lagrange}{Lagrange}

% Boxed final results
\newcommand{\boxedid}[1]{\boxed{\ #1\ }}

% ----------------------------------------------------------------------------
% Title
% ----------------------------------------------------------------------------
\title{The Akita Bridge:\\
\large Binius$_{2^{128}}$-IOP $\,\to\,$ Akita lattice-PCS over $\F_p$}
\author{Binius64 maintainers}
\date{\today}

\begin{document}
\maketitle

\begin{abstract}
This document specifies the \emph{Akita bridge} component of Binius64.
The bridge translates an \emph{oracle relation} stated over the binary
tower field $\binF$ (the field Binius's interactive oracle protocols
speak natively) into a \emph{polynomial-commitment opening} over a prime
field $\akF$ (the field the Akita lattice PCS speaks natively). Three
opening variants are presented: \textsf{full-open} (non-succinct;
testing-only), \textsf{succinct} (succinct, two-point batched opening),
and \textsf{claim-reduced} (succinct, single-point opening via a
claim-reduction sumcheck). Per-step soundness is analysed; the dominant
error term is identified.

This document is the canonical mathematical specification of the
bridge. The Rust implementation in
\verb|crates/akita-bridge/| and
\verb|crates/akita-bridge-prover/| follows section-by-section.
\end{abstract}

\tableofcontents

\bigskip\hrule\bigskip


% ============================================================================
\section{Setting}
\label{sec:setting}
% ============================================================================

Binius64's proof system, on each iteration of its interactive oracle
protocol (IOP), reduces a fragment of the circuit's satisfiability claim
to one or more \emph{oracle relations} of the form
\[
  \langle \Message, \Transparent \rangle \;=\; e
  \qquad \text{(in } \binF \text{)},
\]
where:

\begin{itemize}[leftmargin=*, itemsep=2pt]
  \item $\Message \in \binF^{L}$ is the committed oracle (the prover
        has previously sent a commitment to its multilinear extension);
  \item $\Transparent \in \binF^{L}$ is a \emph{transparent} vector,
        known explicitly to both parties (either materialised or
        evaluable via a closure);
  \item $e \in \binF$ is the claimed inner product;
  \item $L = 2^{\ell}$ where $\ell$ is the log-length of the oracle.
\end{itemize}

The verifier must check that this inner product holds, given a
commitment to $\Message$ but without seeing $\Message$ itself.

The native way to discharge this claim in Binius would be to run a
sumcheck protocol over $\binF$ and discharge the residual evaluation
$\Message(r)$ through a $\binF$-flavoured polynomial commitment scheme
(BaseFold, in the default deployment). The Akita bridge instead routes
this claim through a \emph{lattice-based} PCS that operates over a
\emph{prime field} $\akF$, specifically $\akF \;=\;
\mathsf{Prime128Offset275}$, a 128-bit pseudo-Mersenne prime field. The
bridge's job is to translate the claim from $\binF$-language into
$\akF$-language soundly.

\paragraph{Why bridge to a prime field at all?}
The lattice PCS has two appealing properties Binius does not get from
BaseFold:

\begin{enumerate}[leftmargin=*, itemsep=2pt]
  \item \emph{Post-quantum security.} Akita's binding rests on the
        hardness of Short Integer Solutions (SIS) over module lattices,
        which is plausibly post-quantum-secure under standard
        assumptions. BaseFold's binding rests on collision-resistance
        of the underlying hash family.
  \item \emph{Smaller proofs for moderate-size oracles.} For our
        bit-table workloads at $2^{17}$–$2^{20}$ bits committed, the
        Akita opening is $\sim 80$–$100$\,KB, comparable to or smaller
        than the BaseFold equivalent.
\end{enumerate}

The cost is two non-trivial bridges of mathematics: from a binary tower
field into a prime field (Sections~\ref{sec:bit-slicing} and
\ref{sec:parity-bridge}), and from a multi-oracle interactive argument
into the single-polynomial PCS opening that Akita supports
(Sections~\ref{sec:selected-sum}--\ref{sec:closing}).


% ============================================================================
\section{Notation and preliminaries}
\label{sec:notation}
% ============================================================================

\subsection*{Fields}

\begin{itemize}[leftmargin=*, itemsep=2pt]
  \item $\binF$ is the binary tower field of order $2^{128}$ (Binius's
        $\mathsf{BinaryField128bGhash}$). It is the field over which
        Binius oracle relations are stated.
  \item $\akF$ is the Akita prime field of order $p$, where
        $p = 2^{128} + 275$ (specifically: a pseudo-Mersenne prime). It
        is the field over which the Akita PCS commits and opens.
  \item $\GF{2}$ is the prime subfield $\{0, 1\}$. Both $\binF$ and
        $\akF$ contain $\GF{2}$ as a subset (trivially, in the second
        case: under the identification $0 \leftrightarrow 0_{\akF}$ and
        $1 \leftrightarrow 1_{\akF}$).
\end{itemize}

\subsection*{Multilinear extensions}

For a vector $v = (v_0, \ldots, v_{N-1}) \in \F^{N}$ where $N = 2^n$, its
\emph{multilinear extension} (MLE) $\MLE{v} : \F^n \to \F$ is the unique
multilinear polynomial in $n$ variables such that
$\MLE{v}(\mathbf{i}) = v_i$ for every $\mathbf{i} \in \bits^n$ identified
with its binary representation. Explicitly,
\[
  \MLE{v}(\mathbf{x}) \;=\; \sum_{\mathbf{i} \in \bits^n}
    v_i \cdot \eq(\mathbf{i}, \mathbf{x}),
\]
where $\eq(\mathbf{i}, \mathbf{x}) = \prod_{k=0}^{n-1} \bigl(i_k x_k
+ (1 - i_k)(1 - x_k)\bigr)$ is the \emph{equality polynomial} that
evaluates to $1$ when $\mathbf{i} = \mathbf{x}$ and $0$ otherwise (for
$\mathbf{x} \in \bits^n$). For $\mathbf{x} \in \F^n$ arbitrary,
$\eq(\cdot, \mathbf{x})$ acts as a Lagrange interpolation basis.

\subsection*{Inner products}

For $u, v \in \F^N$ we write $\langle u, v \rangle = \sum_{i=0}^{N-1}
u_i v_i \in \F$ for the usual (componentwise) inner product.
Equivalently, in MLE language,
\[
  \langle u, v \rangle \;=\; \sum_{\mathbf{x} \in \bits^n}
    \MLE{u}(\mathbf{x}) \cdot \MLE{v}(\mathbf{x}).
\]

\subsection*{Constants}

\begin{itemize}[leftmargin=*, itemsep=2pt]
  \item $w := 128$ — the number of bits in one $\binF$ element.
  \item $L := 2^{\ell}$ — the number of $\binF$ elements in the oracle;
        $\ell$ is the \emph{log-msg-len} parameter the bridge takes from
        the caller.
  \item $N := L \cdot w = 2^{\ell + 7}$ — the size of the bit-table
        we will commit to (Section~\ref{sec:bit-slicing}).
  \item $n := \ell + 8$ — the number of variables of the OneHotPoly we
        commit (one extra variable comes from the $\mathsf{K}=2$ encoding
        of each bit; Section~\ref{sec:commitment}).
\end{itemize}


% ============================================================================
\section{The oracle relation}
\label{sec:oracle-relation}
% ============================================================================

We restate the relation to be proved precisely.

\begin{definition}[Oracle linear relation]
\label{def:oracle-relation}
Given:
\begin{itemize}[leftmargin=*, itemsep=2pt]
  \item a committed oracle $\Message \in \binF^L$ (the verifier has a
        commitment to $\MLE{\Message}$ but does not see $\Message$);
  \item a transparent vector $\Transparent \in \binF^L$ (both parties
        agree on it, either materialised explicitly or via an MLE
        evaluator closure);
  \item a public claim $e \in \binF$;
\end{itemize}
the \emph{oracle linear relation} is the predicate
\[
  \langle \Message, \Transparent \rangle \;\stackrel{?}{=}\; e
  \qquad \text{(equality in } \binF \text{)}.
\]
\end{definition}

The bridge's task is to convince the verifier that the relation holds,
in three stages that progressively reshape the claim:

\begin{enumerate}[leftmargin=*, itemsep=2pt]
  \item[\textbf{Stage A.}] Bridge $\binF$-arithmetic into
        $\akF$-arithmetic without leaving any soundness slack
        (Section~\ref{sec:bit-slicing} and
        Section~\ref{sec:parity-bridge}).
  \item[\textbf{Stage B.}] Reduce the resulting $\akF$-side claim to
        evaluations of a single committed polynomial $\MLE{\BitTable}$
        at one or two random points
        (Section~\ref{sec:selected-sum} and
        Section~\ref{sec:booleanity}).
  \item[\textbf{Stage C.}] Discharge those evaluation claims by opening
        a single committed OneHotPoly via Akita's PCS
        (Section~\ref{sec:closing} and
        Sections~\ref{sec:variant-succinct}--\ref{sec:variant-claim-reduced}).
\end{enumerate}

\subsection{What the bridge inherits from Binius (and what it does not)}
\label{sec:inherits-from-binius}

It is worth being explicit about the granularity of the bridge's
input. Binius's IOP layer reduces a circuit's AND / MUL / shift
constraints to a small number of oracle inner-product relations of the
form in Definition~\ref{def:oracle-relation}. Those reductions
\emph{do} involve bit-level reasoning internally (the AND reduction
runs a sumcheck over a $\GF{2}$-bilinear product polynomial, for
example), but their \emph{output} is at field-element granularity:

\begin{itemize}[leftmargin=*, itemsep=2pt]
  \item the committed oracle $\Message$ is a vector of \emph{$\binF$
        elements}, not bits — Binius commits at field-element
        granularity because per-bit commitment would inflate proof and
        witness sizes by a factor of $w = 128$;
  \item the transparent vector $\Transparent$ and the public claim $e$
        are likewise $\binF$ objects;
  \item the only operations Binius itself performs on these objects
        before handing them to the bridge are $\binF$-additions and
        $\binF$-multiplications by public constants.
\end{itemize}

In particular, the prover does \emph{not} hand the bridge a separate
bit-table commitment. The bit-slicing of Section~\ref{sec:bit-slicing}
is the bridge's own re-interpretation of $\Message \in \binF^L$ as a
flat array of $N = L w$ bits, applied freshly inside the bridge — on
the prover side by decomposing each $\Message_i$ into its
$\GF{2}$-basis coordinates, on the verifier side by the same trivial
extraction once the bit-table is committed.

This decomposition is what makes the field-bridging trick of
Section~\ref{sec:parity-bridge} possible: $\binF$-multiplication has no
useful homomorphism into $\akF$, but $\GF{2}$-addition (XOR of bits)
coincides with $\akF$-addition modulo $2$, which is checkable in $\akF$
without slack.

\begin{remark}[Interchangeability]
The bridge is agnostic to how the input relation was produced. Anything
upstream that exposes an oracle linear relation over $\binF$ — Binius64
in our deployment, but also Iron Spartan or a future system — can plug
into the bridge unchanged.
\end{remark}


% ============================================================================
\section{Bit-slicing the oracle}
\label{sec:bit-slicing}
% ============================================================================

The first move is purely structural: we re-express the oracle
$\Message \in \binF^L$ as a flat \emph{bit-table} of $N = L \cdot w
= 2^{\ell + 7}$ Boolean values, so that every subsequent step can
operate on individual bits rather than $\binF$ elements.

\subsection{The \texorpdfstring{$\GF{2}$}{F\_2}-basis decomposition}

$\binF = \GF{2^{128}}$ is a $w$-dimensional vector space over its prime
subfield $\GF{2}$. Fix once and for all an arbitrary $\GF{2}$-basis
$\mathcal{B} = (e_0, e_1, \ldots, e_{w-1})$ of $\binF$.\footnote{In the
implementation we use the implicit basis induced by the
\texttt{BinaryField128bGhash} raw byte layout: each element is stored
as a \texttt{u128}, and bit $j$ of that \texttt{u128} is its
$e_j$-coordinate. The bridge's soundness does not depend on the
specific choice of $\mathcal{B}$ — only that the map
$\binF \to \bits^{w}$ it induces is bijective.} Every $u \in \binF$
has a unique decomposition
\[
  u \;=\; \sum_{j=0}^{w-1} u^{(j)} \cdot e_j,
  \qquad u^{(j)} \in \GF{2}.
\]
We refer to $u^{(j)}$ as \emph{the $j$-th bit of $u$} and denote the
extraction by $\beta_j : \binF \to \GF{2},\; u \mapsto u^{(j)}$.

\subsection{The bit-table}

\begin{definition}[Bit-table]
\label{def:bit-table}
Given the oracle $\Message = (\Message_0, \ldots, \Message_{L-1}) \in
\binF^L$, the \emph{bit-table} is the array
\[
  \BitTable \in \bits^{L \times w},
  \qquad
  \BitTable_{i,j} \;:=\; \beta_j(\Message_i),
  \quad
  0 \le i < L,\ 0 \le j < w.
\]
\end{definition}

\begin{remark}[Layout]
\label{rem:bit-table-layout}
Flattened, $\BitTable$ is a length-$N$ vector with the convention that
\emph{the bit index $j$ varies fastest}. That is, if we define
$\mathsf{flat}(B) \in \bits^N$ by $\mathsf{flat}(B)[i \cdot w + j] :=
\BitTable_{i,j}$, then traversing $\mathsf{flat}(B)$ in order visits all
$w$ bits of element $0$, then all $w$ bits of element $1$, etc. With
$N = 2^{\ell + 7}$, the binary index $\mathbf{x} \in \bits^{\ell + 7}$
of a flat slot splits as
\[
  \mathbf{x} \;=\; \bigl(\underbrace{x_0, \ldots, x_6}_{=\, j \text{ in
  binary}},\; \underbrace{x_7, \ldots, x_{\ell + 6}}_{=\, i \text{ in
  binary}}\bigr),
\]
so the lowest 7 variables address \emph{bit-within-element} and the
remaining $\ell$ variables address \emph{element index}.
\end{remark}

\subsection{Embedding into \texorpdfstring{$\akF$}{F\_p}}

To take the bit-table into Akita-land we embed $\GF{2}$ into the prime
field via the obvious map
\[
  \iota : \GF{2} \hookrightarrow \akF,
  \qquad 0 \mapsto 0_{\akF}, \quad 1 \mapsto 1_{\akF}.
\]
This is a ring homomorphism on the source's additive structure (since
$1_{\akF} + 1_{\akF} = 2_{\akF} \ne 0_{\akF}$, it is \emph{not} a ring
homomorphism overall, but addition modulo $2$ on $\GF{2}$ becomes
ordinary integer addition modulo $p$ on $\akF$, which we will exploit
deliberately in Section~\ref{sec:parity-bridge}). We abuse notation and
write $\BitTable \in \akF^N$ for the image under $\iota$ of the flat
bit-table; this is the object the rest of the protocol manipulates.

\begin{definition}[Bit-table MLE]
\label{def:bit-table-mle}
$\MLE{\BitTable} : \akF^{\ell + 7} \to \akF$ is the unique multilinear
extension of the flat bit-table (post-embedding):
\[
  \MLE{\BitTable}(\mathbf{x})
  \;=\; \sum_{\mathbf{i} \in \bits^{\ell + 7}}
        \iota(\BitTable_{\mathsf{unflatten}(\mathbf{i})})
        \cdot \eq(\mathbf{i}, \mathbf{x}).
\]
By construction $\MLE{\BitTable}(\mathbf{i}) \in \{0_{\akF}, 1_{\akF}\}$
for every $\mathbf{i} \in \bits^{\ell + 7}$ — a fact the verifier will
later enforce \emph{cryptographically} (Section~\ref{sec:booleanity})
because the prover cannot otherwise be trusted to commit to a Boolean
table.
\end{definition}

\subsection{Properties}

The decomposition has three properties we will rely on repeatedly:

\begin{lemma}[Inversion]
\label{lem:bit-table-inversion}
Given the bit-table $\BitTable$, the oracle is recoverable via
$\Message_i = \sum_{j=0}^{w-1} \BitTable_{i,j} \cdot e_j \in \binF$.
The map $\Message \leftrightarrow \BitTable$ is bijective.
\end{lemma}
\begin{proof}
Immediate from uniqueness of $\GF{2}$-basis decomposition.
\end{proof}

\begin{lemma}[Cardinality and dimension]
\label{lem:bit-table-card}
Committing to the bit-table requires $N = L \cdot w = 2^{\ell + 7}$
Boolean values, equivalently $\ell + 7$ MLE variables (one extra
variable in the OneHotPoly encoding — see
Section~\ref{sec:commitment} — brings the variable count to $n = \ell
+ 8$).
\end{lemma}

\begin{lemma}[Linearity of inner products]
\label{lem:inner-product-bilinearity}
For any $u, v \in \binF$, $\langle u, v \rangle_{\binF}$ (the product
in $\binF$) decomposes \emph{linearly} into a $\GF{2}$-bilinear
combination of bits of $u$ and bits of $v$:
\[
  uv \;=\; \sum_{j, k} \beta_j(u) \beta_k(v) \cdot (e_j e_k)
  \quad \text{in } \binF,
\]
and each output bit $\beta_m(uv) \in \GF{2}$ is a fixed $\GF{2}$-linear
combination of the products $\beta_j(u) \beta_k(v)$. The combination
coefficients are determined entirely by the basis $\mathcal{B}$ and are
publicly computable.
\end{lemma}

Lemma~\ref{lem:inner-product-bilinearity} is what makes the parity
bridge of Section~\ref{sec:parity-bridge} possible: each output bit of
$\langle \Message, \Transparent \rangle_{\binF}$ is a parity sum (i.e.\
$\GF{2}$-linear combination) over fixed subsets of the bit-table, so a
verifier holding only the bit-table can independently re-compute every
output bit.

\paragraph{Implementation.}
\verb|BitSliceOracle::from_binius_oracle| in
\verb|crates/akita-bridge/src/protocol.rs| materialises the bit-table
into a $w \times L$ matrix \verb|bit_polys|, with the convention
\verb|bit_polys[j][i] = beta_j(M_i)| (note: bit index first, element
index second — the transpose of Definition~\ref{def:bit-table}). The
flattening of Remark~\ref{rem:bit-table-layout} is realised by
\verb|to_bit_table_evals|, which iterates the element index in the
outer loop and the bit index in the inner loop, recovering the
$j$-varies-fastest ordering.


% ============================================================================
\section{Bridging fields: the parity identity}
\label{sec:parity-bridge}
% ============================================================================

\textsf{[Section to be filled in.]}

\paragraph{Outline.}
Each output bit of $\langle \Message, \Transparent \rangle$ is, viewed as
a $\GF{2}$-linear function of the input bits, a parity sum over a
specific subset of the bit-table. We lift these parity sums from
$\GF{2}$ to integers ($\mathbb{Z}$), bound each integer by the public
table $\Transparent$, and reinterpret modulo $p$ to land in $\akF$. The
verifier then re-checks
\(
   \sum_{i,j} \Transparent^{(j)}_i \cdot \BitTable_{i,j}
   = \text{(opened sum)} \pmod p
\)
as an $\akF$-arithmetic statement on the bit-table.

Soundness: dominated by the probability that the opened sums match
$\langle \Message, \Transparent \rangle$'s parity output bit-by-bit but
diverge in integer-lifted form, which we'll bound by $\negl(\lambda)$ via
the bounded-parity argument.

Implementation pointer: \verb|protocol::BatchedParityBridgeProof|,
\verb|protocol::parity_sum_bounds|.


% ============================================================================
\section{Selected-sum sumcheck}
\label{sec:selected-sum}
% ============================================================================

\textsf{[Section to be filled in.]}

\paragraph{Outline.}
The verifier samples $\alpha \in \akF^*$, defines the
\emph{selected-bit mask}
$\Lambda_\alpha : \akF^{\ell + 7} \to \akF$ as an MLE of a specific
public table built from $\Transparent$ and $\alpha$, and the prover runs
a degree-2 sumcheck on
\(
   \sum_{\mathbf{x} \in \bits^{\ell+7}} \MLE{\BitTable}(\mathbf{x}) \cdot
   \Lambda_\alpha(\mathbf{x}) \,=\, \Sigma_\alpha,
\)
where $\Sigma_\alpha$ is the public $\alpha$-RLC of the opened parity
sums. Output: a random point $\selectedPoint \in \akF^{\ell + 7}$ and an
evaluation claim $\MLE{\BitTable}(\selectedPoint) = \evalSelected$
(communicated by the prover and checked against the sumcheck's closing
claim).

Soundness: $\le (\ell + 7) \cdot 2 / |\akF^*|$ from the sumcheck rounds,
plus $1/|\akF^*|$ from sampling $\alpha$.

Implementation pointer: \verb|protocol::prove_product_sumcheck_transcript|,
\verb|protocol::verify_product_sumcheck_transcript|.


% ============================================================================
\section{Weighted booleanity sumcheck}
\label{sec:booleanity}
% ============================================================================

\textsf{[Section to be filled in.]}

\paragraph{Outline.}
Selected-sum alone does not pin $\MLE{\BitTable}$ to Boolean values
(a soundness gap: a cheating prover could substitute a non-Boolean
function that happens to inner-product to the right number with
$\Lambda_\alpha$). We close the gap with a degree-3 sumcheck
\[
   \sum_{\mathbf{x} \in \bits^{\ell+7}} w(\mathbf{x}) \cdot
   \MLE{\BitTable}(\mathbf{x}) \cdot
   \bigl(\MLE{\BitTable}(\mathbf{x}) - 1\bigr) \;=\; 0,
\]
where $w(\mathbf{x}) = \eq(\mathbf{r}_w, \mathbf{x})$ is the equality
weight at a fresh challenge $\mathbf{r}_w \in \akF^{\ell + 7}$.
Booleanity at every point would force the inner factor to vanish
pointwise; the weighted sum at a random $\mathbf{r}_w$ is sufficient
because $w$ is generic.

Output: another random point $\boolPoint \in \akF^{\ell + 7}$ and
evaluation claim $\MLE{\BitTable}(\boolPoint) = \evalBool$.

Soundness: $\le (\ell + 7) \cdot 3 / |\akF^*|$ from the sumcheck rounds.

Implementation pointer:
\verb|protocol::prove_weighted_booleanity_sumcheck_transcript|.


% ============================================================================
\section{Closing the bridge: PCS opening}
\label{sec:closing}
% ============================================================================

\textsf{[Section to be filled in.]}

\paragraph{Outline.}
After Sections~\ref{sec:selected-sum} and \ref{sec:booleanity}, the
verifier holds two unverified evaluation claims:
\(
   \MLE{\BitTable}(\selectedPoint) = \evalSelected
\)
and
\(
   \MLE{\BitTable}(\boolPoint) = \evalBool.
\)
Both claims must be \emph{discharged} via the PCS opening of the
committed polynomial. Two strategies follow in
Section~\ref{sec:variant-succinct} (two-point batched opening) and
Section~\ref{sec:variant-claim-reduced} (single-point opening after a
claim-reduction sumcheck).


% ============================================================================
\section{Variant A — Succinct (multi-point) opening}
\label{sec:variant-succinct}
% ============================================================================

\textsf{[Section to be filled in.]}

\paragraph{Outline.}
We open the committed polynomial at both $(\selectedPoint \| 1)$ and
$(\boolPoint \| 1)$ simultaneously via Akita's multi-point
\verb|batched_prove| / \verb|batched_verify| API. The trailing $1$
recovers the bit-table MLE from the OneHotPoly encoding
(Section~\ref{sec:commitment}).

Implementation pointer: \verb|akita-bridge/src/succinct.rs|.


% ============================================================================
\section{Variant B — Claim-reduced (single-point) opening}
\label{sec:variant-claim-reduced}
% ============================================================================

\textsf{[Section to be filled in.]}

\paragraph{Outline.}
Sample $\beta \in \akF^*$, run a degree-2 sumcheck on
\[
   \sum_{\mathbf{x}} \MLE{\BitTable}(\mathbf{x}) \cdot
   \bigl(\eq(\selectedPoint, \mathbf{x}) + \beta \cdot
        \eq(\boolPoint, \mathbf{x})\bigr)
   \;=\; \evalSelected + \beta \cdot \evalBool.
\]
The sumcheck output is a single new evaluation claim
$\MLE{\BitTable}(\reducedPoint) = \evalSelected'$ (computed from the
final-claim equation), which we discharge with a single-point Akita
opening at $(\reducedPoint \| 1)$.

Trade-off: one extra degree-2 sumcheck (cheap) in exchange for a
single-point PCS opening (smaller and faster than the two-point
variant).

Implementation pointer: \verb|akita-bridge/src/claim_reduced.rs|.


% ============================================================================
\section{Commitment representation}
\label{sec:commitment}
% ============================================================================

\textsf{[Section to be filled in.]}

\paragraph{Outline.}
We commit to the bit-table as a \emph{2-way OneHotPoly}:
$\mathsf{OneHotPoly}\langle \akF, D=64, K=2 \rangle$. Each Boolean
$b \in \bits$ is encoded as the pair $[1 - b, b]$ (exactly one $1$ in
each pair), padding the bit-table from $N$ to $2N$ field elements and
adding one variable. The resulting polynomial has $n = \ell + 8$
variables. To recover $\MLE{\BitTable}$ at a point $\mathbf{x} \in
\akF^{\ell + 7}$, one opens the OneHotPoly at $(\mathbf{x}\,\|\,1)$:
the trailing $1$ selects the second element of each $[1-b, b]$ pair.

Discussion of why $K=2$ (vs $K=64$, vs DensePoly): in
\verb|docs/akita-k2-onehot-bug-report.md| and Section~\ref{sec:closing}
above.

Implementation pointer: \verb|protocol::BitSliceOracle::to_onehot_bit_table_poly|.


% ============================================================================
\section{End-to-end soundness}
\label{sec:soundness}
% ============================================================================

\textsf{[Section to be filled in.]}

The protocol's soundness error is a sum of per-component contributions:

\begin{table}[h]
\centering
\begin{tabular}{lll}
\toprule
Component & Source of error & Bound \\
\midrule
Bit-slicing                 & --- (perfect)                            & $0$ \\
Parity bridge               & Integer-lift collision                   & TBD \\
Selected-sum sumcheck       & Schwartz–Zippel ($\alpha$ + $\ell + 7$ rounds) & $\le (\ell + 8) \cdot 2 / |\akF^*|$ \\
Booleanity sumcheck         & Schwartz–Zippel ($\mathbf{r}_w$ + $\ell + 7$ rounds) & $\le (\ell + 8) \cdot 3 / |\akF^*|$ \\
Claim-reduction sumcheck (B)& Schwartz–Zippel ($\beta$ + $\ell + 7$ rounds) & $\le (\ell + 8) \cdot 2 / |\akF^*|$ \\
PCS opening (Akita)         & Lattice SIS hardness                     & $\negl(\lambda)$ \\
\bottomrule
\end{tabular}
\end{table}

Dominant contribution: TBD.


% ============================================================================
\appendix

\section{Round-by-round protocol scripts}
\label{app:scripts}
% ============================================================================

\textsf{[Appendix to be filled in.]}

Will contain side-by-side prover/verifier scripts for both variants,
formatted as algorithm2e environments.


% ============================================================================
\section{Implementation correspondence}
\label{app:impl}
% ============================================================================

A short mapping from sections to source files and key types in the
implementation.

\begin{table}[h]
\small
\begin{tabular}{lll}
\toprule
Section & Crate :: module                                & Notable items \\
\midrule
\ref{sec:bit-slicing}        & \verb|akita-bridge::protocol|     & \verb|BitSliceOracle|, \verb|to_bit_table_evals| \\
\ref{sec:parity-bridge}      & \verb|akita-bridge::protocol|     & \verb|BatchedParityBridgeProof|, \verb|parity_sum_bounds| \\
\ref{sec:selected-sum}       & \verb|akita-bridge::protocol|     & \verb|prove_product_sumcheck_transcript| \\
\ref{sec:booleanity}         & \verb|akita-bridge::protocol|     & \verb|prove_weighted_booleanity_sumcheck_transcript| \\
\ref{sec:closing}            & --- (verifier-side check)         & --- \\
\ref{sec:variant-succinct}   & \verb|akita-bridge::succinct|     & \verb|AkitaSuccinctVerifierChannel| \\
                             & \verb|akita-bridge-prover::succinct|  & \verb|AkitaSuccinctProverChannel| \\
\ref{sec:variant-claim-reduced} & \verb|akita-bridge::claim_reduced| & \verb|AkitaClaimReducedVerifierChannel| \\
                             & \verb|akita-bridge-prover::claim_reduced| & \verb|AkitaClaimReducedProverChannel| \\
\ref{sec:commitment}         & \verb|akita-bridge::protocol|     & \verb|BitSliceOracle::to_onehot_bit_table_poly| \\
\bottomrule
\end{tabular}
\end{table}


\end{document}
